% ============================================================================
% NEW MATERIAL, October 5, 2026 (agent draft for Tommaso; not in JMP_DS_draft).
% A complete Model section written to the model as implemented at the base
% point promoted Oct 6 (output/model/production/2007: Estate A + net-estate
% floor, balanced property-tax rebate, PSID levels entry; loss 15.07,
% provisional, native postcheck passed). First drafted on chain 11,
% engine code/model/experiments/birth_count_choice/model/engine/.
% It does not replace sections/03_model.tex, which is left unchanged; the
% October 5 build (JMP_DS_mock_oct05.tex) inputs this file instead.
% Where the code and the author's stated intent differ, the text follows the
% intent and a TODO marks the difference. Calibrated values are placeholders
% until the Estate A refit; fixed inputs carry their source in a comment.
% ============================================================================
\providecommand{\mocktodo}[1]{\textcolor{red}{[TODO: #1]}}

\section{Model}
\label{sec:model-oct05}

Motivated by these observations, I build a life-cycle model in which the same
households choose when to have children, how much housing to occupy, and
whether to rent or own it. Households live in an overlapping-generations
economy, face persistent earnings risk, and save in a single bond. Children
raise the household's needs for consumption and, above all, for space. Space
can be rented or bought. Renting and owning the same dwelling cost the same per
period, but the largest dwellings are available only to owners, owners enjoy
their home somewhat more, selling a home is costly, and buying one requires a
down payment that must be accumulated by saving. House prices and rents are
determined in equilibrium by an upward-sloping supply of housing. The
government runs a pay-as-you-go pension system and taxes property.

\subsection{The economy}

Time is discrete and a period lasts four years.
% [MODEL INPUT] period_years 4.0 (bundle).
I describe the households, their preferences, the earnings process, housing,
the financial market and the government in turn.

\subsubsection*{Demographics.}
A household enters the economy at age 18 and lives at most through age 85, for
$J$ periods indexed by $j$. It works until retirement at age 66 and then
receives a pension. From retirement on, it survives from one period to the
next with probability $s_j<1$; before retirement $s_j=1$.
% [MODEL INPUT] ages 18,22,...,82 (17 cells), J_R = cell 12 (age 66),
% survival 0.9391, 0.9185, 0.8850, 0.8300 for cells 12-15, certain death after
% age 82. Source of the survival schedule not documented in the repo.
\mocktodo{name the source of the survival probabilities (not documented in the
repository; ask Tommaso).}
Each household is one decision unit; I do not model the formation or
dissolution of couples.

Let $n$ denote the number of children ever born to the household and $m$ the
number currently living at home. Births are possible at ages 18 to 45, at most
one per period. Each child at home leaves independently with probability
$\lambda$ per period, so that a child stays on average $1/\lambda$ periods;
departure lowers $m$ but not $n$. I keep track of $n$ up to three, with the
last state standing for three or more.
% [MODEL INPUT] fertile cells ages 18-42 (decision at the start of each
% four-year period, so births up to age 45); lambda = 2/9 per four-year
% period, mean stay 18 years; no newborn exemption; top state n = 3 = "3+".

\subsubsection*{Preferences.}
A household values consumption $c$, housing services $s$, and children at
home. Its flow utility is
\begin{equation}
 u(c,s;m)
 =\frac{1}{1-\sigma}\left(\frac{c^{\alpha}s^{1-\alpha}}{e(m)}\right)^{1-\sigma}
 +\psi\, m^{1-\gamma}.
 \label{eq:m5-utility}
\end{equation}
The first term is utility from the consumption--housing composite per
equivalent adult. The equivalence scale
\begin{equation}
 e(m)=\left(\frac{2+\omega m}{2}\right)^{\nu}
 \label{eq:m5-escale}
\end{equation}
measures how many adult equivalents the household contains, so that the same
spending buys less utility per person in a larger family. The second term is
the direct benefit of children at home, with $\psi>0$ its level and
$\gamma\in[0,1)$ the rate at which the benefit of an additional child
declines.
% [MODEL INPUT] e(m) code: ((2+0.7 m)/2)**0.7, multiplied into utility as
% e**(sigma-1); omega = nu = 0.7, m top-coded at 3. Code form
% e^(sigma-1) * x^(1-sigma)/(1-sigma) equals (x/e)^(1-sigma)/(1-sigma).
\mocktodo{cite the source of the equivalence scale ($\omega=\nu=0.7$); the
September deck cites Scholz et al.\ (2006), not checked against the paper.}

Children need room. A household with children at home obtains services only
from the space in excess of a minimum $h_P$:
\begin{equation}
 s=
 \begin{cases}
 h^R-h_P\,\mathbf 1\{m\geq1\}, & \text{renter of } h^R \text{ rooms},\\[2pt]
 \chi\bigl(h-h_P\,\mathbf 1\{m\geq1\}\bigr), & \text{owner of } h \text{ rooms},
 \end{cases}
 \qquad \chi\geq1 .
 \label{eq:m5-services}
\end{equation}
The parent floor $h_P$ is the space a family needs before additional rooms
add to its welfare. The floor is measured in rooms, and the ownership
premium $\chi$ applies only to the rooms above it. It does not depend on the number of children, so it is
the first child that creates the need for space. The parameter $\chi$
captures the additional services an owner obtains from a dwelling of a given
size, for example from the ability to adapt it.
% [MODEL INPUT/CODE] hbar_child_rooms = 0; owner composite is chi*(h - h_P)
% (kernels.py 858-871), not chi*h - h_P as on the September slides.
% Settled: the floor does not grow with children (open object, fine for now).
% Author decision Oct 6: keep the code form chi*(h - h_P); an experiment with
% chi*h - h_P may follow.
Together, \eqref{eq:m5-utility} and \eqref{eq:m5-services} define what a
child costs. At given prices, a first child raises the spending needed to
reach a given utility in two ways: through the equivalence scale, which
scales up all spending, and through the floor, which must be paid at the
rental price of space before any housing yields services.

\subsubsection*{Fertility.}
At the start of each fertile period, a household chooses whether to try for a
child. An attempt succeeds with probability $\pi_j$, which declines with age.
A first birth carries a one-time utility cost $\xi$ of becoming a parent. The
decision is subject to type-I extreme-value taste shocks with scale $\kappa_1$
when the household is childless and $\kappa_C$ when it already has children.
After the outcome of the attempt is known, the household chooses its housing
and saving for the period.
% [MODEL INPUT] pi_j = 1 - 0.02 exp(0.134 (age - 18)), zero from 45;
% fecundity constants unfitted (open item). Fertility logit scales
% kappa_fert (n = 0) and kappa_fert_continuation (n >= 1).
Housing can therefore adjust to a birth within the same period, which matches
the timing of the room response in Section~\ref{sec:empirical-evidence}.

\subsubsection*{Earnings.}
Before retirement, a household of age $j$ earns
\begin{equation}
 y_j(z)=(1-\tau^w)\,w\,\varepsilon_j\,z ,
 \label{eq:m5-earnings}
\end{equation}
where $w$ is the wage, normalized to one, $\varepsilon_j$ a deterministic
age profile, $\tau^w$ the payroll tax, and $z$ an idiosyncratic productivity
state. Productivity follows a first-order autoregressive process in logs,
\begin{equation}
 \log z'=(1-\rho)\mu_z+\rho\log z+\eta',\qquad \eta'\sim N(0,\sigma_\eta^2),
 \label{eq:m5-ar1}
\end{equation}
where $\mu_z$ normalizes mean productivity to one. Earnings risk is the
only source of income uncertainty. Retirees receive a
pension $\varpi$ that is the same for all retirees.
% [MODEL INPUT] rho = 0.735, sigma_eta = 0.484 per four-year period (annual
% 0.926, 0.270), sigma_z = 0.713, mu_z = -0.252 (E[z] = 1). Method of moments,
% PSID gross labor earnings EARNINDRRC, annual waves 1984-1997, non-overlapping
% four-year block means residualized on age and year, rho = 0.3735/0.5085;
% bootstrap rho [0.692, 0.769]. Traced by the utility thread, Oct 6.
% Earlier note: rho = 0.7346, sigma_eta = 0.4838 at four-year frequency,
% estimated by method of moments on PSID four-year log-earnings
% autocovariances (single_process_external_estimate.json); nine-state
% discretization; age profile from breakpoints 22/26/34/46/58. Pension flat
% (retirement_income_z_scale = 0), level balanced endogenously.
\mocktodo{name the data source of the age profile $\varepsilon_j$ (not
documented in the files read).}

\subsubsection*{Housing.}
Housing is measured in rooms and trades at price $P$ per room. Owners choose
among dwellings of $h\in\mathcal H=\{2,4,6,8,10\}$ rooms. Renters choose any
size up to $\bar h^R$ rooms, so the largest dwellings can only be owned, as in
the American Housing Survey facts of Section~\ref{sec:empirical-evidence}.
Owners pay depreciation $\delta$ and property tax $\tau^H$ on the value of
their home each period. Selling a home costs a fraction $\varsigma$ of its value;
buying costs nothing beyond the price.
% [MODEL INPUT] H = {2,4,6,8,10}; rental cap 6 rooms (retained provisional);
% delta = 1.416% a year, tau_H = 1.060% a year (adopted national inputs);
% selling cost 6%.
A household's tenure and dwelling choice is subject to type-I extreme-value
taste shocks with scale $\kappa_H$, one for each option: renting, keeping the
current home, or buying each size in $\mathcal H$.

Competitive landlords own the rental stock. A landlord buys a room at $P_t$,
pays depreciation and property tax, rents it out at $r_t$, and sells it at
$P_{t+1}$. Landlords finance themselves at the household bond return $R$, so
zero profits require
\begin{equation}
 r_t=(R+\delta+\tau^H)P_t-P_{t+1},
 \qquad\text{and in a steady state}\qquad
 r=(R-1+\delta+\tau^H)P .
 \label{eq:m5-rent}
\end{equation}
Rent per room equals the owner's user cost.
Because the rent equals the user cost and households and landlords borrow and
save at the same rate, a room costs the same per period whether it is rented or
owned. Tenure matters for the cost of a family through four features only: the
size limit on rentals $\bar h^R$, the ownership premium $\chi$, the cost of
selling $\varsigma$, and the down payment described next.

Housing is supplied by a competitive construction sector with a constant
elasticity with respect to the rental price,
\begin{equation}
 H^s_t=H_0\left(\frac{r_t}{\bar r}\right)^{\zeta},
 \label{eq:m5-supply}
\end{equation}
where $H_0$ sets the scale of the stock, $\bar r$ is a normalizing rent and
$\zeta$ is the supply elasticity.
% [MODEL INPUT] xi_supply = 0.63 (pinned; the 1.75 eta_supply field is unread);
% r_bar = 0.16; H0 = 6.2996 derived at N = 1. Whether supply should respond to
% gross- or net-of-tax rent is an open item in CALIBRATION_STATUS.md.

\subsubsection*{Saving and mortgages.}
Households save in a one-period bond with gross return $R$, fixed in the world
market. Let $b$ denote the bond position a household carries into a period;
$b<0$ is mortgage debt. Borrowing must be secured by a home: renters cannot
borrow, and an owner of $h$ rooms can borrow up to a fraction $\phi$ of the
home's value,
\begin{equation}
 b'\geq-\phi\,P h .
 \label{eq:m5-ltv}
\end{equation}
A buyer must therefore end the period of purchase holding equity of at least
$(1-\phi)Ph$, the down payment. Since $b'$ is the outcome of the household's
consumption--saving choice, the down payment is accumulated through saving;
there is no separate rule for how it is financed. The constraint applies to
any new borrowing. An owner whose debt exceeds $\phi Ph$ after a fall in the
price is not forced to repay, but cannot borrow more. Mortgages carry no
amortization schedule and no limit on payments relative to income.
% [CODE] native_due_stayer_credit: stayer floor min(b, -phi P h);
% mortgage_amortization 0; no PTI. phi = 0.80 (fixed, non-negotiable).
% The engine also screens purchases on R b + y; that screen is implied by
% (eq:m5-ltv) at the end of the period and is not written here (see the
% October 1 assessment in memory).

\subsubsection*{Budget constraints.}
Transactions in housing take place after interest accrues on the bond carried
into the period. Let $S=(1-\varsigma)P h_{-}$ denote the proceeds from selling the
home a household enters the period with, if any, and $Q=Ph$ the price of a home
it buys. A household that rents $h^R$ rooms faces
\begin{equation}
 c+r\,h^R+b'=Rb+y+T+S,
 \qquad b'\geq0,
 \label{eq:m5-bc-rent}
\end{equation}
and a household that owns $h$ rooms at the end of the period faces
\begin{equation}
 c+(\delta+\tau^H)P h+b'=Rb+y+T+S-Q,
 \qquad b'\geq-\phi Ph,
 \label{eq:m5-bc-own}
\end{equation}
with $S=0$ for a household that did not own and $S=Q=0$ for an owner who
keeps its home. Here $y$ is earnings or the pension and $T$ a lump-sum
transfer from the government. The sale proceeds and the purchase price are not
compounded: a household that sells and buys within a period carries only the
difference between them into next period's interest.
% [CODE] post-interest timing, author-adopted: b' = R b + S - Q + y - c - K
% (kernels.py 301/316/331, 379/394/409; household.py 306). The September
% slides still show the original timing b' = R(b + S - Q) + y - c - K.

\subsubsection*{Estates.}
A household that dies leaves its net estate, the bond plus the value of its
home net of the selling cost, $W=b'+(1-\varsigma)Ph'$. It values the estate
according to
\begin{equation}
 v(W)=\theta_0\,\frac{(\theta_1+W)^{1-\sigma}-\theta_1^{1-\sigma}}{1-\sigma},
 \label{eq:m5-bequest}
\end{equation}
where $\theta_0$ governs the strength of the bequest motive and $\theta_1$ the
extent to which bequests are a luxury. Estates are not received by other
households in the model; entrants' wealth is drawn from the data instead, as
described below.
% [CODE] Estate A: bequest_utility_vec(b + (1-psi) P h), normalized,
% theta_n = 0, estate tax 0, no recipient; settled as intended for now.
% Net-estate floor b' >= -(1-psi) P h' where death is possible; slack at
% phi = 0.8, so not stated.

\subsubsection*{Government.}
The government runs two programs, each with a balanced budget. The pension
system pays every retiree the same pension $\varpi$, financed by the payroll
tax $\tau^w$ on earnings. The property tax raises $\tau^H$ times the value of
the entire housing stock, owned and rented, and its revenue is returned to
all households as an equal lump-sum transfer $T$.
% [MODEL, base_2007, Oct 6 refit 15.07] balanced rebate on inside every solve
% (tmp/rebate_entry_impl_20261006; output/model/production/2007). Chains 11
% and 13 had no rebate; that discrepancy is resolved at the new base.
% [MODEL INPUT] pension_mode balanced_stationary; tau_w = 0.0803 from the CPS
% ASEC pension/earnings ratio 0.2294; pension 0.918 per period (derived).
% The September slides cite a 17.9% payroll tax; the code uses 0.0803.

\subsubsection*{Entry.}
Households enter at age 18 as childless renters. Each entrant draws its
productivity from the stationary distribution of $z$. Its wealth comes from the
distribution of net worth among childless renters aged 18 to 24 who head their
household in the PSID, 1984 to 2019, with negative net worth set to zero.
Dollar amounts are converted to model units with one common scale, chosen so
that mean income in this sample equals mean earnings of entrants in the model,
and observations are assigned to productivity states by income rank, so that
the poorest entrants in the data become the least productive entrants in the
model. The number of entrants in each period equals the number of children
born sixteen to twenty years earlier divided by $2.1$, so that the population
is constant when each household has $2.1$ children on average.
% [MODEL, base_2007] entry law levels_rank (author decision Oct 5-6):
% closure_root/.../model/entry_wealth.py; N = 1,835 family-years, 29.6% with
% negative net worth (censored at zero). Earnings in the scale exclude all
% transfers. Replaces the chain-11 five-bin rescaled rule.
% [CODE] adult_entry.py: entry = adjusted births / 2.1, half entering 16
% years after birth and half 20; adjusted births count a birth into the
% top state as 3.602 children (provisional). The 0.5 entrant_conversion_factor
% is read only by an outside-flow valve that is off at chain 11.

\subsection{The household's problem}
\label{sec:m5-household}

The state of a household at the start of a period is its age $j$,
productivity $z$, bond position $b$, number of children ever born $n$ and at
home $m$, and housing $h_{-}$, which is either a home it owns or zero for a
renter. A period has three stages. The household first chooses whether to try
for a child, then chooses tenure and dwelling, and finally chooses consumption
and saving. Children leave home and productivity is drawn between periods.

\subsubsection*{Saving.}
Consider a household that has decided to own $h$ rooms and enters the saving
stage with cash on hand $x=Rb+y_j(z)+T+S-Q$. Its value is
\begin{multline}
 V^O_j(x,z,n,m,h)=\max_{c,\,b'}\;u\bigl(c,\chi(h-h_P\mathbf 1\{m\geq1\});m\bigr)\\
 +\beta\,\mathbb E\Bigl[s_j\,\mathcal W_{j+1}(b',z',n,m',h)+(1-s_j)\,v\bigl(b'+(1-\varsigma)Ph\bigr)\Bigr]
 \label{eq:m5-vown}
\end{multline}
subject to \eqref{eq:m5-bc-own}. The expectation is over next period's
productivity $z'$ and the number of children $m'$ still at home. The renter's
value $V^R_j(x,z,n,m)$ is defined in the same way, with the choice of $h^R\leq
\bar h^R$, the renter's services, and the constraint \eqref{eq:m5-bc-rent}.
Because the composite in \eqref{eq:m5-utility} is Cobb--Douglas, a renter not
at the size limit spends a share $1-\alpha$ of its outlay net of the floor on
rooms above $h_P$.

\subsubsection*{Tenure and dwelling.}
Let $o$ index the options available to a household entering the period with
housing $h_{-}$: renting, keeping its home if it owns one, or buying a dwelling
of each size in $\mathcal H$. Each option implies its own cash on hand and its
own saving-stage value, $V_o$. An option is available only if it leaves room
for a saving choice that satisfies the borrowing constraint. With the
extreme-value shocks, the value of entering this stage is
\begin{equation}
 \mathcal H_j(b,z,n,m,h_{-})=\kappa_H\log\sum_{o}\exp\bigl(V_o/\kappa_H\bigr),
 \label{eq:m5-tenure}
\end{equation}
and the probability of choosing option $o$ is
$\exp(V_o/\kappa_H)/\sum_{o'}\exp(V_{o'}/\kappa_H)$.

\subsubsection*{Fertility.}
At the start of a fertile period, a household with $n<3$ compares not trying,
which leaves it at $\mathcal H_j(b,z,n,m,h_-)$, with trying, which yields
\[
 U^1_j=\pi_j\bigl[\mathcal H_j(b,z,n+1,m+1,h_-)-\xi\,\mathbf 1\{n=0\}\bigr]
 +(1-\pi_j)\,\mathcal H_j(b,z,n,m,h_-).
\]
With scale $\kappa=\kappa_1$ if $n=0$ and $\kappa=\kappa_C$ otherwise, the
value at the start of the period is
\begin{equation}
 \mathcal W_j(b,z,n,m,h_-)=\kappa\log\Bigl[\exp\bigl(\mathcal H_j(b,z,n,m,h_-)/\kappa\bigr)
 +\exp\bigl(U^1_j/\kappa\bigr)\Bigr].
 \label{eq:m5-fertility}
\end{equation}
Outside the fertile ages, or with $n=3$, $\mathcal W_j=\mathcal H_j$.

\subsection{Equilibrium}

I study stationary equilibria, in which prices and the distribution of
households over states are constant, and transitions between them.

\begin{definition}
A stationary equilibrium consists of a house price $P$ and rent $r$; a pension
$\varpi$ and payroll tax $\tau^w$; a transfer $T$; decision rules for
fertility, tenure, dwelling size, consumption and saving; and a distribution
$\mu$ of households over states, such that
\begin{enumerate}
\item the decision rules solve the household problem of Section~\ref{sec:m5-household} given
prices, taxes and transfers;
\item the rent satisfies the landlord's zero-profit condition
\eqref{eq:m5-rent};
\item the housing market clears: the rooms occupied by owners and renters
equal the supply \eqref{eq:m5-supply},
\[
 \int\Bigl[\mathbf 1\{\text{own}\}\,h+\mathbf 1\{\text{rent}\}\,h^R\Bigr]\dd\mu
 =H_0\left(\frac{r}{\bar r}\right)^{\zeta};
\]
\item the pension system balances, $\tau^w w\int\varepsilon_j z\,
\mathbf 1\{j<J_R\}\dd\mu=\varpi\int\mathbf 1\{j\geq J_R\}\dd\mu$, and the
transfer returns property-tax revenue, $T\int\dd\mu=\tau^H P\int\bigl[
\mathbf 1\{\text{own}\}\,h+\mathbf 1\{\text{rent}\}\,h^R\bigr]\dd\mu$;
\item $\mu$ is the stationary distribution implied by the decision rules, the
earnings process, survival, the departure of children, and entry, with the
mass of entrants equal to past births divided by $2.1$.
\end{enumerate}
\end{definition}

Two features of the equilibrium are worth stating. First, the price of
housing is what links fertility to housing supply: more births raise the
demand for rooms by families, which raises $P$ and $r$ and with them the cost
of the space a child requires. Second, because the population is constant
only if each household has $2.1$ children on average, the scale of housing
supply $H_0$ and the level of fertility are tied together in a stationary
equilibrium. In the quantification I normalize the population to one and
choose $H_0$ so that the market clears at the price consistent with a constant
population.%
\footnote{Equivalently, the price is the one at which the stationary
population reproduces itself, and $H_0$ is the supply scale that makes that
price clear the housing market. A counterfactual holds $H_0$ fixed and lets
the population adjust.}
\mocktodo{check this closure paragraph against Tommaso's intended
wording; the code solves the price from the renewal condition and derives
$H_0$ (production/README.md lines 18--19).}

% ----------------------------------------------------------------------------
% ----------------------------------------------------------------------------
% Parameter values for Section 4 must be read from the base point:
%   python3 code/model/tools/read_results.py fit base_2007
%   output/model/production/2007/parameters_estate_a.csv
% (Oct 6 refit, Mac chain 7, case 0129_nm, loss 15.0712, provisional;
% e.g. beta 0.9397, psi 0.157, chi 1.081, h_P 2.585, H0 6.312). Not typeset.
% ----------------------------------------------------------------------------

\paragraph{Open items for the author (mock only).}
\mocktodo{(1) Whether the income
gradient of Table~\ref{tab:emp-fertility-income} is an untargeted evaluation
fact, and whether to recompute it on the 2004 and 2006 CPS. (2) The wording of
the equilibrium closure (price from the renewal condition, $H_0$ derived).
(3) Sources for the survival probabilities, the age profile $\varepsilon_j$
and the equivalence scale.}
